% Source: https://artofproblemsolving.com/community/c3h2057757_2018_aime_i_1
% Stable identity: aops-topic-2057757-q01
A factorisation is $x^2+ax+b=(x+p)(x+q)$ with integers $p,q$, so $p+q=a$ and $pq=b$. Since $a\ge1$ and $b=pq\ge0$, both $p,q$ are non-negative (they cannot both be negative as $p+q>0$, and if one were negative and the other positive then $pq<0$). So for each $a$, the possible $b$ are the values $p(a-p)$ for $p=0,1,\dots,a$, where $p$ and $a-p$ give the same value. The number of distinct values is the number of $p$ with $0\le p\le\dfrac {a}{2}$, i.e.\ $\left\lfloor\dfrac {a}{2}\right\rfloor+1$ (these values $p(a-p)$ are strictly increasing in $p$ on this range).

Hence
\[
S=\sum_{a=1}^{100}\left(\left\lfloor\dfrac {a}{2}\right\rfloor+1\right)=100+\sum_{a=1}^{100}\left\lfloor\dfrac {a}{2}\right\rfloor=100+2(1+2+\dots+49)+50=100+2450+50=2600,
\]
and the remainder is $\boxed{600}$.
